\section{Introduction} \label{sec:introduction}
Foams are among the most familiar yet physically complex systems encountered in everyday life. From the froth on your morning coffee to the structure of biological tissues, foams occupy a unique place in the study of soft matter. At their core, a foam is a dispersion of bubbles seperated by a liquid medium, and the collective behaviour of these bubbles; how they coarsen over time and how the foam restructures because of this, is they core of what is being investigated in this report.\\

Coaresening is a fundamental property of these foams. It is the process by which bubbles exchange gas across their shared interfaces causing larger bubbles to get larger and smaller bubbles get smaller. Over time, the average bubble area increases and the number of bubbles decrease. During this coarsening, bubbles gain new neighbours as the smaller ones disappear causing topological rearrangements in the foam. Understanding the statistical character of the individual bubble trajectories.